\documentclass[11pt]{article}

% ---------- packages ----------
\usepackage[a4paper,margin=1in]{geometry}
\usepackage{amsmath,amssymb,amsthm}
\usepackage{booktabs}
\usepackage{longtable}
\usepackage{tabularx}
\usepackage{microtype}
\usepackage[hidelinks]{hyperref}
\usepackage{enumitem}
\usepackage{xcolor}
\usepackage{titlesec}

\setlist{itemsep=2pt,topsep=2pt}
\titleformat{\section}{\large\bfseries}{\thesection}{0.6em}{}
\titleformat{\subsection}{\normalsize\bfseries}{\thesubsection}{0.6em}{}

% ---------- shorthand ----------
\newcommand{\Lvec}{\mathbf{L}}
\newcommand{\LJ}{\Lvec_{J2000}}
\newcommand{\Iten}{\mathbf{I}}
\newcommand{\J}{\mathbf{J}}
\newcommand{\wvec}{\boldsymbol{\omega}}
\newcommand{\qvec}{\mathbf{q}}
\newcommand{\R}{\mathbf{R}}
\newcommand{\NLL}{f_{\mathrm{NLL}}}
\newcommand{\ttt}[1]{\texttt{#1}}
\newcommand{\eq}[1]{Eq.~(\ref{#1})}
\newcommand{\rf}[1]{RF~Eq.~(\ref{#1})}  % Robinson--Frueh equation

% ---------- metadata ----------
\title{s073e --- Deep audit of Robinson \& Frueh 2025\\
       (J.~Astronaut.~Sci.~73:7, DOI 10.1007/s40295-025-00557-9)\\
       \large Part 1: implementation lessons we can adopt.\\
              Part 2: does their work encroach on our s073 framing?}
\author{Girish (PI) --- with Claude as scribe}
\date{2026-05-14}

\begin{document}
\maketitle

\begin{abstract}
\noindent
Robinson \& Frueh 2025 (RF25 hereafter) is the highest--remaining--threat paper
flagged by the s073c literature--novelty audit: a 37--page treatment of
global light--curve attitude inversion under measurement noise and inertia
uncertainty, by the closest competitor group (Frueh, Purdue). This document is a
deep read of RF25 with two narrow goals.
\textbf{Part 1} catalogues every implementation element of RF25 that could
materially improve our pipeline, with concrete decisions on whether to adopt,
adapt, or pass. \textbf{Part 2} asks whether their work encroaches on the
specific framing s073 is hypothesising --- that the body--frame Casimirs
$(|\Lvec|, 2T)$ plus polhode phase parameterise an inertial--$\LJ$ ambiguity that
is generic to asymmetric--inertia bodies. The conclusion of Part~2: the framing
space is \emph{not} closed by RF25, but they have observed adjacent phenomena
(continuous $\wvec$--bands on axisymmetric bodies) that we must explicitly
distinguish from our cat--4 claim if we want to publish.
\end{abstract}

\section*{Reading guide}
\noindent
Treat this as a working document, not a paper. It is more thorough than a
reading note because we may end up citing or reproducing some of these claims;
it is less polished than a manuscript section because the cat--4 framing is
still a hypothesis (s073d is ambiguous as of commit \ttt{b6e3482}). Numbers and
quotes come from the PDF at \ttt{/home/girish/Documents/s40295-025-00557-9.pdf};
specific page citations use the \emph{print} page numbers shown at the top of
each PDF page (e.g.\ ``p.~5627'' = the journal's page 5627, not the PDF's nth
sheet). Citations to RF25's own bibliography use their numbering ([RF\#nn]).

% =====================================================================
\section{What RF25 is, in one paragraph, before we judge it}
% =====================================================================

\noindent
RF25 builds a global light--curve attitude--inversion pipeline that explicitly
optimises through (a) realistic per--timestep measurement noise and (b)
uncertainty in the object's principal moments of inertia. The state vector is
$\mathbf{x} = (p_1, p_2, p_3, \omega_1, \omega_2, \omega_3, J_y/J_x,
J_z/J_x)^T$ --- three MRP (modified Rodriguez parameter) attitude components,
three body--frame angular velocity components, and two inertia--tensor ratios.
The loss is a negative--log--likelihood under a per--timestep Gaussian model.
The solver is parallel BFGS started from $10^5$ uniform initial guesses
(uniform on $SO(3)$, uniform in an $\omega$--ball, Gaussian on inertia ratios).
The forward model uses a Blinn--Phong BRDF and a semi--analytic Sutherland--Hodgman
polygon--clipping shadow algorithm precomputed into a 4D body--frame
``shadow cache.'' They demonstrate on two objects: an ATLAS V rocket body (1987--084D,
axisymmetric inertia) and a HYLAS--4 box--wing satellite (2014--004A,
asymmetric inertia). They run with both low ($\sigma_J = 0.1$) and high
($\sigma_J = 1.0$) inertia--ratio prior uncertainty. All cases use a model
mismatch between the truth shape and the assumed inversion shape, deliberately
introduced to simulate real--data conditions.

% =====================================================================
\part{Implementation lessons we can adopt}
\label{part:implementation}
% =====================================================================

\noindent
Twelve elements of RF25 stand out as candidates for transplantation into our
pipeline. We sort them by ``how much they would change our code'' and tag each
with a \textbf{decision} of \textsc{adopt} (drop in directly), \textsc{adapt}
(useful idea, our regime differs), \textsc{evaluate} (worth a small experiment
to test), or \textsc{pass} (their constraint, not ours).

\subsection{The negative--log--likelihood (NLL) loss with per--timestep~$\sigma_k$ weighting and signal--rescaling \textnormal{[\textsc{adopt}, high--value]}}

\noindent
RF25's loss function is their single most exportable contribution to our work.
Their derivation (their \S2.1, p.~5--6 of the article):

\begin{equation}
\label{eq:rf_nll}
\NLL(\hat{\mathbf{S}}, \mathbf{S})
= \frac{1}{m} \sum_{k=1}^{m}
\left[
\ln \sigma_k \;+\; \tfrac{1}{2} \left( \tfrac{S_k - \hat{S}_k}{\sigma_k} \right)^2
\right]
\end{equation}

\noindent
where $\hat{S}_k$ is the predicted ADU flux at timestep $k$, $S_k$ is the
measured ADU flux, and $\sigma_k$ is the predicted measurement standard
deviation (computed from the full noise model in their \S3.5.1; see Item~\ref{ssec:noise}).
The $\ln \sigma_k$ term penalises solutions that ``predict noisy regions
confidently''; the squared term is the usual chi--square.

\textbf{Why we should care.} Our current cost surfaces use surrogate MSE on
magnitude. RF25 argue (their p.~4--5) that magnitude--space $\ell_2$ ``attributes
much larger errors to brighter regions of the light curve, leading to solutions
that overfit bright specular glints while ignoring smaller variations in the
signal,'' and is blind to per--timestep noise. \emph{This is precisely the bright--peak
overfit pattern we have seen} --- e.g.\ s058 / s059 LM polish on a 21--epoch
window converging to phantom basins disagreeing with the full--LC hi--fi by
$\rho = 30$--$70$ (see memory \ttt{project\_local\_window\_polish\_phantom\_basins.md}).

A subtle but load--bearing trick they add is the \textbf{signal rescaling}: before
computing the loss, $\hat{\mathbf{S}}$ is multiplied by
$\|\mathbf{S}\| / \|\hat{\mathbf{S}}\|$, forcing the prediction and observation
to have equal $\ell_2$ norm in flux (their p.~5).
This absorbs systematic albedo--area errors and prevents the loss from
penalising LCs of the correct shape but wrong scale. \textit{This is essentially
a free correction for the kind of BRDF coefficient uncertainty that haunts our
m048 forward model.}

\textbf{Concrete adoption path.} Replace the surrogate--MSE residual in
\ttt{lib/forward.py} (and in the LM cost in \ttt{s058/s059} polish stages)
with $\NLL$ on ADU flux, computed with our existing per--timestep $\sigma_k$
(\ttt{lib.noise} already returns this for the m048 noise model). Add the
$\|\mathbf{S}\|/\|\hat{\mathbf{S}}\|$ rescaling factor inside the residual.
This is a $\sim$ 30--line change. No retraining of the surrogate is needed
because the surrogate maps body--frame coordinates $\to$ magnitude, not
loss--space; we just re--score on top of the same evaluator.

\textbf{Caveat.} Our $\rho$--band convention is rooted in surrogate MSE
ratios, not NLL. If we adopt NLL we either (a) report both during a
transition, or (b) re--derive the band thresholds. Either is cheap.

\subsection{The full sensor + sky noise model (Eqs.~35--40) \textnormal{[\textsc{evaluate}, scoped to real--data work]}}
\label{ssec:noise}

\noindent
RF25 enumerate \emph{four independent} noise contributors and combine them per
timestep (their \S3.5.1, p.~17--19):

\begin{itemize}
\item \textbf{Background Poisson noise} (their Eq.~37):
$\lambda_{\mathrm{back},k} = n_{\mathrm{pix}} \cdot
(\lambda_{\mathrm{moon}} + \lambda_{\mathrm{star}} + \lambda_{\mathrm{twi}}
+ \lambda_{\mathrm{zod}} + \lambda_{\mathrm{air}} + \lambda_{\mathrm{poll}})$.
This is scattered moonlight, integrated background starlight, twilight,
zodiacal light, airglow, and light pollution, each as an independent
Poisson term.
\item \textbf{Scintillation noise} via Young's atmospheric approximation
(their Eq.~38): a function of telescope aperture diameter $D$, exposure
time $\Delta t_k$, observation zenith angle $\gamma_k$, and observing
station altitude $h_{\mathrm{obs}}$.
\item \textbf{Sensor noise} (their Eq.~35): dark current (Poisson, rate
$\lambda_{\mathrm{dark}}$ per pixel per second) plus read noise
(Gaussian, $\sigma_{\mathrm{read}}$). The flat--field term scales as the
square root of total signal across pixels (their Eq.~36).
\item \textbf{Shot noise} on the object signal (their Eq.~39):
$\lambda_{\mathrm{shot},k} = \bar{S}_{\mathrm{SO},k} / \sqrt{g}$.
\end{itemize}

\noindent
Total per--timestep variance is their Eq.~(40):
$\sigma_{S,k}^2 = \lambda_{\mathrm{back},k}
+ \bar{S}_{\mathrm{SO},k} \sigma_{Y,k}^2
+ \lambda_{\mathrm{shot},k} + \sigma_{\mathrm{flat}}^2 + \sigma_{\mathrm{sensor}}$.

\textbf{Why we should care.} Our current noise model uses a fixed apparent--magnitude
noise floor (the standard $m_{\mathrm{noise}} = 0.05$ mag). RF25 demonstrate
(their Fig.~6c, 12c) that the noise envelope varies by 4--5$\times$ across a
single light curve: faint epochs are scintillation--dominated, bright epochs
are read--noise dominated. \emph{This is precisely the regime where the
$\NLL$ loss buys you something over surrogate MSE.}

\textbf{Decision.} \textsc{evaluate} not \textsc{adopt}, because our entire
benchmark cohort (m048 seeds 0--119) was generated with the simpler noise
model. Swapping to the full noise model mid--cohort would invalidate the
post--fix replications (s068--s071) we just committed. The right path is:

\begin{enumerate}
\item Implement the RF25 noise model in a new module \ttt{lib/noise\_rf.py}
that returns $\sigma_k$ as a function of (geometry, telescope params,
flux level). Wire it as an optional flag.
\item Generate a \emph{small} replication cohort (10 seeds, RF--noise) and
verify the surrogate still resolves Band A--C basins. The surrogate is
bridge--independent (memory \ttt{feedback\_surrogate\_is\_bridge\_independent.md})
so it should not need retraining.
\item If basin structure changes materially, the full noise model becomes a
secondary publication thread; if not, fold it in for the real--data
chapter only.
\end{enumerate}

\subsection{The shadow cache \textnormal{[\textsc{evaluate}, high--leverage if forward--model wall is the bottleneck]}}
\label{ssec:shadow_cache}

\noindent
RF25 build a 4D boolean lookup table of shape $n_{\mathrm{Az}} \times
n_{\mathrm{El}} \times n_f \times n_f$ indexed by (azimuth, elevation, shadower
face, shadowed face) in the \emph{body frame}. Their \rf{eq:rf_cache}:

\begin{equation}
\label{eq:rf_cache}
\hat{\mathbf{e}} = \begin{bmatrix} \cos(\mathrm{El}) \cos(\mathrm{Az}) \\
                                   \cos(\mathrm{El}) \sin(\mathrm{Az}) \\
                                   \sin(\mathrm{El}) \end{bmatrix}
\end{equation}

\noindent
At evaluation time, the body--frame Sun and observer directions are converted
to (Az, El), rounded down (their floor--operator scheme, Eqs.~28--29) to the
nearest cache entry, and the per--face shadowing is read directly. At
$n_{\mathrm{Az}} = n_{\mathrm{El}} = 200$ and $n_f = 100$, the cache is
\textbf{50 MB} (one bit per entry). They report the truncation error to
nearest cached cell is ``negligible'' because the only faces that flip due to
the rounding are those that are nearly tangent to the illuminator.

\textbf{Why we should care.} The shadow computation is currently a notable
fraction of our forward--model wall time. For Path~2 (s072 closed--form $q(t)$)
in particular, where the dynamics step is now microseconds, shadow ray--casting
becomes the dominant cost. The s073b cprofile shows the $\phi$--ODE is 81\% of
Path 2 wall (\ttt{experiments/s073b\_path2\_cprofile.md}); after we replace
that with elliprj, shadow cost rises in the ranking.

\textbf{Decision.} \textsc{evaluate}. Build a one--object prototype on m048
geometry, measure speedup vs the trimesh batched ray cast. If the speedup is
$\geq$~5$\times$ on a representative LC, integrate; if less, the engineering
cost (cache invalidation when components articulate) is not worth it.

\textbf{Caveat for articulated geometry.} m048 has solar panels that sun--track
(\ttt{src/articulation/behaviors/sun\_tracking.py}). RF25's cache assumes
\emph{static} geometry in the body frame. We would need either (a) separate
caches per panel orientation, or (b) compute static--component shadows from the
cache and articulated--component shadows from ray casting. Option (b) matches
how our current shadow engine already factors the mesh (cf.\ \ttt{src/computation/shadow\_engine.py}'s
static/articulated separation), so adoption is mechanical.

\subsection{Semi--analytic Sutherland--Hodgman shadowing \textnormal{[\textsc{pass} for now]}}

\noindent
The clipping algorithm RF25 use to compute the shadowed area on each face
(their \S3.3.1--3.3.3, p.~11--14) replaces pixel--wise ray casting with
analytic polygon clipping. Eq.~24 is the inclusion--exclusion formula for
mutual higher--order shadow intersections:

\begin{equation}
\bar{a} = a - \sum_{d=1}^{n} \sum_{K \in \mathrm{comb}(n,d)}
(-1)^{d+1} A\!\left( \bigcap_{k \in K} \mathcal{P}_k \right)
\end{equation}

\noindent
truncated at depth $d=2$ in their work. They report ``nearly identical'' light
curves to pixel--wise reference (their Fig.~6d and 12d).

\textbf{Decision.} \textsc{pass}. We already use trimesh's batched ray
intersection, which is well--optimised and handles arbitrary non--convex
geometry without the truncation--depth trade--off. The Sutherland--Hodgman
approach is most useful when the geometry is exactly polygonal (which RF25's
simplified models are); our m048 STL is a 30k+ triangle mesh where the
inclusion--exclusion combinatorics grow uncomfortably. The shadow--cache
(Item~\ref{ssec:shadow_cache}) likely subsumes the speedup that would
motivate switching the underlying algorithm.

\subsection{State parameterisation: MRP attitude + inertia ratios \textnormal{[\textsc{evaluate} for asymmetric--inertia uncertainty work; \textsc{pass} for our main pipeline]}}

\noindent
RF25 search over $(p_1, p_2, p_3, \omega_1, \omega_2, \omega_3, J_y/J_x, J_z/J_x)$,
with $J_x \equiv 1$. Two pieces here.

\textbf{(a) MRP vs quaternion.} The MRP \rf{eq:rf_mrp_to_dcm} maps to the
DCM via $C(\mathbf{p}) = I_3 + (8\tilde{p}\tilde{p} - 4(1 - \|\mathbf{p}\|^2)\tilde{p})
/ (1 + \|\mathbf{p}\|^2)^2$ (their Eq.~9). MRPs avoid the
unit--norm constraint of quaternions and are 3--dimensional, simplifying BFGS.
They have a singularity at the antipode of identity which RF25 handle via
shadow MRPs (their reference [37]).

\textbf{Decision: pass.} Our quaternion infrastructure is mature, the
propagator--fix work hardened it, and the convention bug experience
makes us cautious about changing representations. The MRP advantage is mainly
for unconstrained optimisation, which is not our bottleneck.

\textbf{(b) Inertia ratios in the search state.} This \emph{is} interesting.
RF25 directly optimise $J_y/J_x$ and $J_z/J_x$ jointly with $(q, \omega)$.
For real--data work on objects without precise inertia knowledge, this is
the right architecture. Their Case 2a (box--wing, low uncertainty) achieves
0.78\% median $|\omega(0)|$ error; Case 2b (high uncertainty) achieves 2.36\%.
\textit{The inertia uncertainty does not prevent inversion convergence on
the asymmetric body.}

\textbf{Decision: \textsc{evaluate}.} For the m048 cohort (known inertia by
construction) this is unnecessary. But the Frueh group's framing is
``real--data inversion needs inertia--ratio search'' and ours is implicitly
``we know the inertia exactly.'' If we want a real--data chapter we will
need this. Slot it for after the cat--4 framing is settled.

\subsection{Loss--based convergence classification: the $\alpha$ threshold \textnormal{[\textsc{adapt}]}}

\noindent
RF25 define $f_{\mathrm{best}}$ as the loss obtained by placing the truth state
in the assumed--model forward simulation (their \S2.5.1, p.~8). Then:
solutions with $f_i < f_{\mathrm{best}}$ are flagged as \emph{overfits}
(discarded); solutions with $f_{\mathrm{best}} < f_i < \alpha \cdot
f_{\mathrm{best}}$ are accepted, with $\alpha \in [1.5, 2]$ chosen by the
user.

\textbf{Why we should care.} This is structurally the same construct as our
$\rho$--band convention but with two important differences:
\begin{itemize}
\item It explicitly \emph{filters out overfits} (any solution below the
truth--model loss). We do not do this. We have spent significant
debugging time on solutions that scored ``better than truth'' because of
surrogate--vs--hi--fi gap or noise overfitting.
\item It is parameterised by an interpretable $\alpha$ that the user can
loosen if the model mismatch is high (RF25 used $\alpha = 2$ for the
matched--shape cases and $\alpha = 1.5$ for the box--wing cases with
heavier mismatch).
\end{itemize}

\textbf{Decision: adapt.} Define an analogous threshold in $\rho$--space:
$\rho_{\mathrm{best}} := \rho(\hat{S}(\text{truth}), S_{\mathrm{obs}})$,
discard $\rho_i < \rho_{\mathrm{best}}$ as overfits, accept up to
$\rho_i < \alpha \cdot \rho_{\mathrm{best}}$ for some $\alpha$ calibrated to
the m048 noise level. This would have caught the s058 oracle--injection
artefact earlier.

\subsection{Initial--condition sampling: uniform $SO(3)$ + $r^{1/3}$--scaled ball for $\omega$ \textnormal{[\textsc{evaluate}]}}

\noindent
RF25's sampling (their \S2.3, p.~6):
\begin{itemize}
\item \textbf{Attitude:} sample $(g_1, g_2, g_3, g_4) \sim \mathcal{N}(0,1)$
independently, normalise to unit length; this gives a uniform quaternion
on $S^3$ (Marsaglia 1972). Then convert to MRP via \rf{eq:rf_q_to_mrp}
$\mathbf{p} = \mathbf{q}_{1:3} / (1 + q_4)$. The uniform--$SO(3)$ property
holds for our quaternion grid too if we sample $S^3$ uniformly.
\item \textbf{Angular velocity:} sample $r \sim \mathcal{U}(0,1)$,
scale by $r^{1/3}$, multiply by $\omega_{\max}$, and direct uniformly on
$S^2$ (again Marsaglia). The $r^{1/3}$ factor gives uniform volume
density inside the 3--ball (their Eq.~8).
\end{itemize}

\textbf{Why we should care.} Our current $\omega$ sampling is a Fibonacci
direction grid $\times$ magnitude grid (\ttt{project\_omega\_grid\_architecture.md}).
This is structured, deterministic, and aligned with the polhode geometry, but
\emph{biases magnitude} (uniform over the magnitude grid is not uniform over
the ball). RF25's uniform--ball sampling is unbiased.

\textbf{Decision: \textsc{evaluate}.} For surrogate--scoring (where the cost
is per--sample) Sobol or Fibonacci with explicit magnitude weighting is fine.
For BFGS or LM polish starts (where the cost is per--converged--point),
uniform--ball is more honest. Add as an option, A/B--test on a 10--seed cohort
slice.

\subsection{The BFGS solver with $10^5$ parallel starts in Taichi \textnormal{[\textsc{pass} for now, \textsc{evaluate} for GPU port]}}

\noindent
RF25 use BFGS from $10^5$ initial conditions, compiled in Taichi for parallel
execution (their \S2.4--2.5, p.~7). The full Case 1a took 9.8~minutes on an
Apple M1 CPU. Comparable: our s059k with $N_{\mathrm{DIRS}}=800$ Fibonacci grid
$\times$ 50 polish starts is 105--127~min per seed on Pool(24+8) (memory
\ttt{project\_s059k\_cohort\_viable\_band\_A.md}). The factor--$\sim$10
speedup they report is real, but most of it is from (a) the simpler forward
model and (b) Taichi's GPU--compiled inner loop.

\textbf{Decision.} \textsc{pass} on BFGS itself --- our LM polish is well--tuned
and the BFGS Hessian approximation is not better in this regime. \textsc{evaluate}
the Taichi port if the m048 forward model becomes the rate--limiter again.

\subsection{Inertia ratios visualisation \textnormal{[\textsc{adopt} as analysis tool]}}

\noindent
RF25 plot the converged $(J_y/J_x, J_z/J_x)$ histogram (their Fig.~8b, 10b,
13b, 14b). For the axisymmetric rocket body the inertia--ratio cloud spreads
along the $J_z/J_x$ axis (Case 1a, Fig.~8b); for the asymmetric box--wing it
collapses onto truth (Case 2a, Fig.~13b). This is a clean diagnostic for
``does the LC discriminate the inertia tensor.''

\textbf{Decision: adopt} as a plot we should add to the s073d / s074
diagnostic kit. If our hypothesised cat--4 family at fixed $(|\Lvec|, 2T)$
exists, it should appear as a particular curve in the $(J_y, J_z)$--ratio
plane (because $|\Lvec|^2 = \sum_i J_i^2 \omega_i^2$ and $2T = \sum_i J_i
\omega_i^2$ couple them).

\subsection{Solution--space visualisation in body--frame ($\omega_x, \omega_y, \omega_z$) \textnormal{[\textsc{adopt}]}}

\noindent
RF25's Fig.~7b is a 3D scatter of converged $\omega$ solutions in body
frame, with the truth in red. The cylindrical bands they describe (axisymmetric
case) are visually obvious in this plot. For the asymmetric case (Fig.~8d,
13d, 14d) they plot $(\omega(0)$ azimuth, $\omega(0)$ elevation) as a 2D
density histogram.

\textbf{Decision: adopt.} We have produced similar plots ad hoc in s053 and
s059i; standardise the plot script so every multi--solution analysis ends with
this view. For cat--4 specifically, the polhode signature (closed curve in
body frame $\omega$) is the right thing to look at, and this is the canonical
view.

\subsection{The ``averaging destroys solution--space structure'' point \textnormal{[\textsc{adopt} as policy]}}

\noindent
RF25's \S4.3, p.~32: ``While it may be attractive to apply an averaging
procedure for the initial orientations or angular velocities to produce a
single state vector, doing so will generally destroy the structure of the
solution space spanned by the original solution set. Instead, we recommend
that densely clustered regions of the solution space may be thinned without
impacting this structure.''

\textbf{Decision: adopt as policy.} This crystallises a discipline we have
been intuitively following but never written down. Add to memory under
\ttt{feedback\_multi\_solution.md}.

\subsection{Ground--truth $f_{\mathrm{best}}$ requires placing truth in the assumed model \textnormal{[\textsc{adopt} as gate]}}

\noindent
Subtle but important. RF25 do not compute $f_{\mathrm{best}}$ from truth in
the truth--model; they compute it from truth in the \emph{assumed} (inversion)
model. The reason: when there is model mismatch, the truth state in the
truth model produces a noiseless LC that, when scored against noisy
observations, yields a residual purely from noise. But the truth state in
the assumed (mismatched) model yields a residual that captures
both noise \emph{and} model--mismatch bias --- which is the right denominator
for ``how good can any solution legally be.''

\textbf{Decision: adopt.} Our ``noise floor'' for $\rho$ is set by the
m048 noise model alone; we should add an analogous quantity that places
truth in a deliberately--perturbed forward model (e.g.\ surrogate v2 vs
hi--fi) and use that as the discard threshold for overfits.

% =====================================================================
\part{Does RF25 encroach on the s073 framing?}
\label{part:encroachment}
% =====================================================================

\noindent
This is the question s073c flagged as load--bearing. RF25 is the closest
competitor paper to the cat--4 framing; if it pre--empts us we redirect.

\subsection{The s073 framing, stated precisely so it can be compared}

\noindent
What s073 measured (verified, source \ttt{notebooks/inversion/survey/results/s073/summary.json}):
on m048 seed~89, the Band A polish output \ttt{cluster\_457} and the truth
attitude have $|\Lvec|$ matching to 0.55\% but $\LJ$--directions differing by
$128.6^\circ$ at $t=0$, with hi--fi $\rho = 0.904$ (within the noise floor).
The s073 framing extrapolates this single observation into a structural claim:

\begin{quote}
\textbf{Hypothesised cat--4 ambiguity:} \emph{For a torque--free rigid body
with no shape symmetry, an LC computed under a fixed inertial geometry
$(\hat{u}, \hat{s})$ over a short observing arc depends on the body--frame
sun/observer trajectory only --- which is determined by $(|\Lvec|, 2T)$
plus a polhode--phase scalar, not by the inertial orientation of $\LJ$. Thus
there exists a continuous 2--parameter family of inertial $\LJ$--directions
producing similar LCs.}
\end{quote}

\noindent
The framing has two distinct claims, which we must judge separately against
RF25:
\begin{description}
\item[Claim A (mechanism)] LC depends on body--frame geometry only; therefore
the inertial direction of $\LJ$ is (to leading order) unobservable on a
short arc.
\item[Claim B (cardinality)] The ambiguity is \emph{continuous} (a
2--parameter manifold over the $S^2$ of $\LJ$--directions, modulo discrete
flip).
\item[Claim C (generality)] This applies on \emph{asymmetric--inertia}
bodies --- it is not the cat--2 ``visual symmetry'' phenomenon.
\end{description}

\noindent
We assess each below.

\subsection{What RF25 actually report about solution--space geometry}

\noindent
RF25 enumerate ambiguities through their empirical solution clouds, not
through theoretical analysis. The relevant passages and figures:

\paragraph{(p.~23, Case 1a, axisymmetric rocket body, $\sigma_J = 0.1$.)}
Quote: ``Figure 7a shows that the initial state MRP solutions lie on
\textbf{two disjoint tubes}. The two families of orientation solutions
correspond to reflecting the assumed cylinder about the phase angle
bisector plane spanned by the vector halfway between the observer and Sun
and their cross product --- \textbf{identical to the mediatrix plane
identified in [33]}.''

This is the cat--1 flip problem (Burton--Robinson--Frueh 2024). The two
``tubes'' are the two sides of the flip mirror.

\paragraph{(p.~23--25, Case 1a continued.)}
Quote: ``Figure 7b reveals that the angular velocity solutions lie mainly
along \textbf{four cylindrical sections symmetric about the body's axis of
symmetry to maintain a constant precession rate}.''

The cylinders are continuous 1D bands in body--frame $\omega$. RF25 attribute
them to the axisymmetric--inertia structure of this rocket body: rotation
about the symmetry axis is dynamically free (the precession rate is set by
$|\omega(0)|J_t / J_a$ from their Eq.~14, while the spin rate
$\omega_s = (J_t - J_a)/J_t \cdot \omega_a(0)$ contributes nothing to the
LC when the symmetry axis is also a geometric axis of symmetry).
\emph{This is the cat--2 ``visual symmetry'' phenomenon, in its axisymmetric
form.}

\paragraph{(p.~25, Case 1a, spin rate observability.)}
Quote: ``The spin rate of this axisymmetric object about its axis of
symmetry is weakly observable in the light curve (at best), we cannot
reliably determine the angular velocity magnitude or the axial moment of
inertia. This observed ambiguity does not contradict well--accepted methods
for angular velocity determination that use brightness ratios [77] or the
periodogram [78], which do not attempt to estimate axial rotation rates
and assume the object's attitude is well--approximated as a single--axis
spin.''

Median spin--rate error: $-95.60\%$ (Table~7) --- basically undetermined.

\paragraph{(p.~27, Case 2a, asymmetric box--wing.)}
Quote: ``Compared to Case 1, there is better convergence in the inertia
tensor ratios as the object has three unique moments of inertia and no
geometric axis of symmetry, \textbf{removing the spin rate ambiguity seen
in the rocket body cases}. While this light curve does not uniquely
determine the initial angular velocity direction, the angular velocity
magnitude estimates have a median close to the truth.''

Median $|\omega(0)|$ error: 0.78\% (Case 2a, Table~10), 2.36\% (Case 2b,
Table~11).

\paragraph{(Figs.~13c--d, 14c--d.)}
For the asymmetric box--wing, the initial orientation--angle error scatters
across $0^\circ$--$180^\circ$ with no clear clustering structure visible at
the resolution of their 2D histograms. The $(\omega(0)$ azimuth, $\omega(0)$
elevation) plots also show scattered solutions across both hemispheres
without obvious continuous bands.

\paragraph{(p.~33, Conclusion, summary of ambiguities.)}
Quote: ``Our results indicate that under realistic conditions of
measurement noise and model mismatch, estimating the orientation of a
space object uniquely from its light curve is unlikely. Even if attitude
inversion indicates only one solution, the non--convex nature of the
problem makes it difficult to determine if one or more better solutions
were missed.''

\subsection{Encroachment analysis, claim by claim}

\begin{table}[h]
\centering
\small
\caption{Each s073 claim versus what RF25 actually states.}
\label{tab:encroachment}
\begin{tabularx}{\linewidth}{@{} l X X l @{}}
\toprule
\textbf{Claim} & \textbf{What RF25 says} & \textbf{Encroachment verdict}
& \textbf{Risk} \\
\midrule

A. \emph{LC depends on body--frame geometry only}
& \emph{Stated in their references} (BRF24~\S4 is cited as the source of the
flip--problem; the body--frame--only dependence is the explicit foundation of
BRF24). RF25 do not re--state it.
& \textbf{Not novel} --- the body--frame--only dependence is the
acknowledged starting point of the Frueh--group programme. Anything we say
about it has to cite BRF24.
& Low. Use as setup, not finding.
\\[6pt]

B. \emph{Continuous 2--parameter $\LJ$--family on asymmetric bodies}
& RF25 demonstrate \emph{continuous 1--parameter bands in body--frame
$\omega$} for the \emph{axisymmetric} case (Fig.~7b, four cylindrical
sections about the body symmetry axis). For the \emph{asymmetric} case
(Fig.~13d, 14d) the solution structure is widely--distributed but not
analysed as a continuous family.
& \textbf{Partially open.} RF25 demonstrate cat--2 (axisymmetric) continuous
$\omega$--bands but explicitly attribute them to body symmetry. They do not
identify a continuous family on asymmetric bodies. They also do not compute
or plot $|\Lvec|$, $2T$, or $\LJ$--direction for any of their solution sets.
& \textbf{This is the gap our framing fills, if the framing converts.}
\\[6pt]

C. \emph{Generic to asymmetric--inertia bodies}
& RF25's asymmetric case (box--wing, Cases 2a/2b) has \emph{many fewer}
converged solutions (90 and 112 vs.\ 2357 for Case 1a) and the authors
\emph{do not} characterise the residual multi--solutions by Casimir
invariants or $\LJ$--direction. They attribute the residual multiplicity
to ``the non--convex nature of the problem'' (p.~33, Conclusion).
& \textbf{Open.} RF25's data may well contain cat--4 instances among
the 90 Case 2a solutions; their analysis would not have surfaced them.
& Medium. If a careful re--analysis of their published results revealed
cat--4 structure their work would still claim priority on the dataset,
even though they did not interpret it that way. \emph{We should not
re--analyse their data; we should build the case on our own m048 results.}
\\

\bottomrule
\end{tabularx}
\end{table}

\subsection{The most dangerous overlap, in detail}

\noindent
The single passage in RF25 that comes closest to our framing is the
description of the \emph{axisymmetric} rocket body Case 1a (p.~23--25):
they observe continuous 1D $\omega$--bands and identify them as preserving
the precession rate while varying the spin rate. The precession rate is
$\omega_p = |\omega(0)| J_t / J_a$ (their Eq.~14). Notice: \emph{this is a
function of $|\Lvec|$ at fixed inertia.} The spin rate $\omega_s$ varies
along the band, and the body--frame polhode (for axisymmetric inertia) is
a circle in the $\omega$--plane perpendicular to the symmetry axis, traced at
the precession rate while the body spins at $\omega_s$ about the symmetry
axis. The cat--2 axisymmetric phenomenon they describe \emph{is literally
``LC depends on $(|\Lvec|, 2T)$ only, modulo spin phase about the symmetry
axis.''}

The danger: a careless reader (or referee) sees our cat--4 framing and
says ``isn't this just RF25's Case 1a generalised?'' We need a clean answer
to that critique, and the answer has to live in two places:

\begin{enumerate}
\item \textbf{For axisymmetric inertia,} body--frame polhode--invariance
\emph{plus visual rotational symmetry of the body} together produce
the continuous spin--rate ambiguity RF25 observes (cat--2 + axisymmetric
inertia). The body--frame--polhode geometry alone cannot explain it
without the visual symmetry --- because for asymmetric body geometry on
the same axisymmetric inertia, the spin rotation would still change the
LC by rotating non--symmetric facets through the body--frame
sun/observer directions.

\item \textbf{For asymmetric inertia + asymmetric body geometry} (our
m048 regime), there is no symmetry to invoke. The body--frame polhode is
not a circle, the rotation has no axis of symmetry, and yet
\emph{cluster\_457 on seed 89 has $|\Lvec|$ matching truth to 0.55\%
but $\LJ$--direction differing by $128.6^\circ$.} This is the empirical
observation that RF25 has not made and that, if it converts to a
class--level finding, fills the gap.
\end{enumerate}

\noindent
The key insight is that the \emph{rotational degree of freedom in inertial
space} (orienting $\LJ$ in $J2000$) is logically independent of any body
symmetry. RF25 collapses it for the axisymmetric body because their
particular case has visual symmetry; on an asymmetric body the same
inertial--rotation DOF should still exist (because the LC depends on
body--frame geometry only) but RF25 do not look for it. Our s073 measurement
is consistent with that DOF being present on asymmetric bodies.

\subsection{What RF25 does \emph{not} provide}

\noindent
A precise inventory of the gaps in RF25 that the s073/cat--4 framing could
fill:

\begin{itemize}
\item \textbf{No use of body--frame Casimir invariants.} ``Polhode'' does
not appear anywhere in RF25. ``Conservation'' is mentioned only in the
context of the BRDF being energy--conserving. The body--frame
$(|\Lvec|, 2T)$ structure is not invoked.

\item \textbf{No computation or plotting of $\LJ$ direction across solutions.}
This is the most striking absence. They have $10^5$ converged solutions
across four cases and never plot the inertial $\LJ$ direction. Our
s073 plot (the one we should produce as a follow--up) is a direct
fill of this gap.

\item \textbf{No analytical proof of the body--frame--only LC dependence
extended to L--direction ambiguity.} BRF24~\S4 proves it for the
discrete flip; the continuous version (claim A above) follows by the
same argument applied to any inertial rotation that preserves
$(\hat{u}, \hat{s})$, but RF25 / BRF24 do not write that argument down.

\item \textbf{No multi--anchor / two--epoch architecture.} RF25 is purely
single--epoch initial--condition optimisation. The s073 idea of using
$\LJ$ conservation at two anchors as a discriminator (s060) is
orthogonal to their approach.

\item \textbf{No connection to the asteroid--LC literature} (Kaasalainen
2001), which is the only place where a continuous--L--family
discussion exists in the literature. They cite the Frueh-group papers
heavily but not Kaasalainen.
\end{itemize}

\subsection{Verdict on encroachment}

\noindent
\textbf{RF25 does not pre--empt the cat--4 framing.} It documents adjacent
phenomena (cat--2 continuous $\omega$--bands on axisymmetric bodies) using
machinery (BFGS clouds, $10^5$ starts) that should have produced the cat--4
observation if it exists, but \emph{they did not analyse their solutions
through the body--frame Casimir lens that would surface it.}

The publication path remains open under two conditions, neither of which is
secured yet:

\begin{enumerate}
\item \textbf{Empirical: cat--4 must be more than N=1.} The cohort sweep
proposed in s073c step~3 (other Band--A multi--sols across seeds 10, 14,
28, 91, \ldots) must show the pattern repeats. If only seed 89's
cluster\_457 exhibits the $(|\Lvec|, 2T)$--match with $\LJ$ mismatch,
the framing is one anomaly.

\item \textbf{Theoretical: the analytical proof of claim A's continuous
version.} A few lines of algebra showing that for any rotation $R \in
SO(3)$ acting on $\LJ$ that preserves the inertial sun and observer
directions over the observing arc, the LC is invariant. The argument is
that the LC depends on body--frame $(\hat{s}_{\mathrm{body}}(t),
\hat{o}_{\mathrm{body}}(t))$, which is determined by the body--frame
trajectory $q_{\mathrm{body}}(t)$ around $\LJ$; rotating $\LJ$ in
inertial space and simultaneously rotating the initial body orientation
by the same $R$ leaves the body--frame trajectory unchanged.
This is the analogue of BRF24's flip argument, written for arbitrary
rotations of $\LJ$ rather than the single $180^\circ$ rotation about
the bisector. It should be a half--page derivation.
\end{enumerate}

\noindent
\textbf{The risk if we ignore RF25.} A referee will ask ``how is this
different from RF25 Case 1a?'' The answer is: ``RF25 Case 1a's continuous
ambiguity is cat--2 (visual + dynamical symmetry); ours is cat--4 (dynamical
only); the discriminator is whether the body shape carries the symmetry,
and our box--wing m048 geometry does not.'' We need to make this distinction
front and centre in any writeup.

\subsection{Robinson \& Frueh 2025 ECSD~2025 conference paper [RF\#74]}

\noindent
RF25 cite their own ECSD 2025 paper ([RF74]: ``Optimal light curve attitude
inversion with measurement noise: two case studies'') as a companion work.
That paper is not in our possession. \emph{Action item:} check ECSD 2025
proceedings for the cited paper; it may add to or weaken the encroachment
assessment, particularly if it explicitly treats asymmetric--inertia
multi--solution structure.

% =====================================================================
\section*{Bottom line, both parts together}
% =====================================================================

\noindent
\textbf{Part 1 (implementation).} The single highest--value adoption is the
NLL loss with $\sigma_k$ weighting and signal rescaling (\eq{eq:rf_nll}). Cheap
to implement, well--motivated, addresses the bright--peak overfitting we have
seen empirically. Second--tier: the shadow cache (if forward--model wall
becomes the bottleneck post--elliprj), the $\alpha$--overfit gate, the inertia
ratio visualisation, and the body--frame $\omega$ scatter plot. Pass on:
MRPs, Sutherland--Hodgman, BFGS swap.

\textbf{Part 2 (encroachment).} RF25 does not close the cat--4 framing. The
adjacent finding they do publish --- continuous $\omega$--bands on
axisymmetric bodies (Case 1a) --- is cat--2 with axisymmetric inertia and a
visually--symmetric body, not the cat--4 hypothesis. The publication path is
still open conditional on (i) showing cat--4 holds beyond N=1 via the
cohort sweep, and (ii) writing the half--page analytical extension of BRF24's
flip argument to continuous rotations of $\LJ$. RF25's [RF74] companion
paper is a residual risk and should be obtained.

\textbf{Next compute, if Part 2 conclusions are accepted.} Two cheap
actions:
\begin{enumerate}
\item Plot $\LJ$--direction for the converged clusters of seed 89 (s069
results), seed 10 (s069), seed 14 (s068), seed 28 (existing surveys) ---
all already in checkpoints. Re--scoring + plotting, no new BFGS, $\leq$30 min
total. This is the cohort sweep step 3 from s073c, on cached data.
\item Write the analytical extension of BRF24 \S4 to continuous $\LJ$
rotation (Part 2 condition ii). No compute, $\sim$an hour of
derivation, can be done in the same session.
\end{enumerate}

\noindent
If both clear, the framing has moved from N=1 anecdote to cohort--validated
class--level claim plus theorem --- and a writeup is in scope.

\end{document}
